\chapter{La ricerca in produzione: quattro difetti di Documentale}\label{chapter:produzione}

I difetti di Koskidex dei capitoli~\ref{chapter:lessicale} e~\ref{chapter:ibrido},
tranne il recupero congiuntivo, erano ipotesi scritte leggendo il codice e poi
misurate. Quelli di questo
capitolo no: sono venuti fuori per strada, lavorando all'innesto e fabbricando
le prime query sul corpus di dominio, e sono difetti della ricerca di Documentale
\emph{in produzione}, cioè di Elasticsearch interrogato come lo interroga l'app.
Sono quattro, e tutti colpiscono le ricerche più ordinarie di un documentale:
un atto di cui si sa il numero, una parola italiana preceduta da un articolo,
un numero insieme al nome del comune, una data. Il quarto è venuto fuori per
ultimo, ed è misurato in un altro modo (sezione~\ref{sec:date}).

Il metodo è lo stesso per tutti. Il difetto si misura su Elasticsearch, passando
dall'app o con la sua query esatta; la causa si isola con un controllo che cambia
una cosa sola; la correzione si misura in Koskidex innestato, dietro una variabile
d'ambiente di Documentale spenta per default. Misurarla in Koskidex ha un senso
preciso grazie alla parità del capitolo~\ref{chapter:sistemi}: partendo dagli
stessi insiemi di Elasticsearch, ogni differenza viene dalla correzione.

\section{Il numero d'atto}\label{sec:numero-atto}

In un documentale si cerca spesso un atto di cui si conosce il numero. Due atti
del Comune di Crispiano fanno da esempio: l'\emph{Ordinanza N.~187} e la
\emph{Determina N.~1223}. La tabella~\ref{tab:numero-atto} dà la posizione
dell'atto giusto con la query di produzione e con due varianti che isolano le
cause.

\begin{table}[ht]
  \centering
  \begin{tabular}{lcc}
    \toprule
    Configurazione & \texttt{ordinanza 187} & \texttt{determina 1223} \\
    \midrule
    Elasticsearch, produzione & \textbf{12° su 85} & \textbf{5° su 8} \\
    Elasticsearch senza refusi & 1° su 1 & 1° su 1 \\
    Elasticsearch senza il campo \texttt{name} & 1° su 84 & 5° su 8 \\
    \bottomrule
  \end{tabular}
  \caption{Posizione dell'atto giusto per due query con il numero d'atto.}
  \label{tab:numero-atto}
\end{table}

\paragraph{La causa sono i refusi sui numeri.} Per \texttt{fuzziness: AUTO}
\texttt{187} è una parola di tre caratteri e ammette un refuso: combacia con
\texttt{18}, \texttt{17}, \texttt{137}, \texttt{186}. Spenti i refusi, l'atto
giusto è primo e unico in tutti e due i casi. Un refuso è una tolleranza pensata
per le parole, dove \emph{manutenzone} è quasi certamente \emph{manutenzione}; per
un numero di protocollo una cifra diversa è un altro atto.

\paragraph{Il nome troncato la aggrava.} Prima di indicizzarlo, l'app passa il nome
del documento da \texttt{pathinfo()} per togliere l'estensione del file. L'atto
che vince \texttt{ordinanza 187} si chiama \emph{ORDINANZA N.~18.2026}, e
\texttt{pathinfo()} prende \texttt{.2026} per un'estensione: nell'indice diventa
\emph{ORDINANZA N.~18}, e nel campo che pesa di più, \texttt{name}, \texttt{18}
combacia con \texttt{187} per un refuso. Per la stessa ragione il nome
dell'atto giusto, dentro l'indice, il numero l'ha perso. Tolto quel campo dalla query, per
l'ordinanza l'atto giusto torna primo; per la determina no, perché lì i
concorrenti hanno il numero vicino nel riassunto.

Due atti sono un esempio, non una misura. La sezione~\ref{sec:numero-comune}
ripete la domanda su trecento atti, e lì il difetto dei refusi sui numeri
ricompare, dietro un altro.

\section{Le elisioni}\label{sec:elisioni}

\paragraph{Il difetto.} Il tokenizer standard di Elasticsearch tiene dentro la
parola un apostrofo fra due lettere (sezione~\ref{sec:differenze}): nell'indice
\emph{dell'illuminazione} è un termine solo, e chi cerca \emph{illuminazione} non
trova un atto che la nomina solo così. L'analizzatore \texttt{italian} di
Elasticsearch ha un filtro apposta, \texttt{elision}, con l'elenco degli articoli
italiani~\cite{es-italian}, ma l'indice di Documentale usa l'analizzatore
predefinito.

\paragraph{Quanti atti tocca.} La misura è fatta sull'indice che l'app ha costruito,
senza ricostruire il tokenizer. Dai termini che Elasticsearch ha scritto davvero
si prendono quelli elisi (lettere, apostrofo, lettere) e per ognuno la parola
dopo l'apostrofo. Una coppia (atto, parola) è \emph{a rischio} se nell'atto la
parola compare solo elisa. È \emph{persa} se la query di produzione con quella sola
parola non trova l'atto: si interroga il motore, invece di contare i termini,
perché i refusi o un plurale nello stesso atto potrebbero salvarla. Il
controllo della causa è un indice temporaneo con lo stesso mapping più il filtro
\texttt{elision}: una coppia che lì torna era persa per l'elisione e non per
altro. Le coppie in cui la forma elisa sta solo nel nome dell'ente
(\emph{Gradisca d'Isonzo}) si contano anche a parte, perché chi cerca un comune ne
scrive probabilmente il nome con l'apostrofo.

\begin{table}[ht]
  \centering
  \begin{tabular}{lrr}
    \toprule
    & atti & \% \\
    \midrule
    nell'indice & 10.018 & \\
    con almeno un termine eliso & 4.280 & 42,7 \\
    \textbf{con almeno una coppia persa} & \textbf{4.033} & \textbf{40,3} \\
    \quad escluso il solo nome dell'ente & 3.826 & 38,2 \\
    \midrule
    & coppie & \\
    \midrule
    a rischio & 5.913 & \\
    salvate dalla ricerca (refusi, forme vicine) & 138 & 2,3 \\
    \textbf{perse} & \textbf{5.775} & \\
    perse che il filtro \texttt{elision} recupera & 5.739 & 99,4 \\
    \bottomrule
  \end{tabular}
  \caption{Le elisioni nell'Elasticsearch di produzione, su 507 parole.}
  \label{tab:elisioni}
\end{table}

\textbf{Quattro atti su dieci hanno almeno una parola che la ricerca di quella
parola non trova} (tabella~\ref{tab:elisioni}). Non dipende dalla fonte: a
Crispiano sono il 34,6\% degli atti, nel Friuli Venezia Giulia il 40,6\%. È la
lingua. La ricerca non salva quasi niente, il 2,3\% delle coppie, e il filtro ne
recupera il 99,4\%: la causa è quella. Delle 36 coppie che restano, 17 sono
elisioni fuori dall'elenco degli articoli (\emph{sant'}, \emph{cinquant'},
\emph{quest'}), 13 sono parole attaccate per uno spazio mancante
(\emph{scuoladell'infanzia}), le altre accenti scritti come apostrofo
(\emph{conformita'}) e prefissi fuori elenco.

Le parole più colpite dicono chi se ne accorge. \emph{Infanzia} è persa in tre atti
su quattro fra quelli che la nominano (168 su 220), perché si scrive quasi sempre
\emph{scuola dell'infanzia} o \emph{nido d'infanzia}; \emph{associazione} in 170 su
230; \emph{art}, cioè i rinvii normativi come \emph{dell'art.~36}, in 595 su 1.406.
\emph{Illuminazione}, la parola da cui tutto è partito, è fra le meno colpite: 12
atti su 99.

\paragraph{La correzione, nei due motori.} In Elasticsearch è il filtro
\texttt{elision} nel mapping e una reindicizzazione. In Koskidex innestato è
un'impostazione, \texttt{ElisionArticles}, con gli stessi 21 articoli
dell'analizzatore \texttt{italian}, che Documentale accende con
\texttt{KOSKIDEX\_ELISION=true}. Misurata con le stesse 507 parole:
\begin{itemize}
  \item senza, Koskidex perde \textbf{le stesse 5.775 coppie} di Elasticsearch in
        produzione, non solo lo stesso numero;
  \item con, ne perde \textbf{le stesse 36} di Elasticsearch col filtro;
  \item sulle 24 query del confronto nessuna query perde un documento, e 8 ne
        guadagnano 65 in tutto: \emph{illuminazione pubblica} passa da 62 a 73
        atti, \emph{contributo associazioni} da 17 a 48;
  \item sulle 300 query known-item della sezione seguente non cambia niente, nemmeno
        l'ordine dei primi dieci.
\end{itemize}
Una differenza fra i due motori resta, ed è a favore di chi cerca: Koskidex toglie
gli accenti, Elasticsearch con l'analizzatore predefinito no. Su 20 delle 507
parole Koskidex trova qualcosa di più, 195 atti in tutto, 67 dei quali per
\emph{unità}: chi scrive \emph{identita} senza accento trova \emph{identità}.

\section{Numero e comune: il campo unico e i refusi sui numeri}\label{sec:numero-comune}

\paragraph{Le query.} Il caso della sezione~\ref{sec:numero-atto} si può
fabbricare in grande senza che nessuno annoti. Da 300 atti estratti a caso con un
seme fisso si scrive la query \texttt{<numero> <comune>}, per esempio
\texttt{13605 Fiume Veneto}; l'atto di partenza è l'unica risposta giusta. Si
tengono solo gli atti la cui coppia (comune, numero) è unica nel corpus, perché
la numerazione riparte ogni anno. Le query e il loro uso sono descritti nel
capitolo~\ref{chapter:valutazione}; qui la metrica principale è l'MRR@10, con la
quota di query con l'atto al primo posto, entro i primi dieci e a vuoto.

\paragraph{Il difetto.} In produzione \textbf{l'atto giusto è entro i primi dieci
nel 2\% delle query}, e sette query su dieci non trovano niente
(tabella~\ref{tab:numero-comune}). Nelle 89 che trovano qualcosa l'atto giusto c'è
in 6. La causa non è il numero ma come la query combina i campi: con
\texttt{best\_fields} e \texttt{operator: and} tutte le parole devono stare nello
stesso campo, ma il numero sta in \texttt{additional\_data} e il comune in
\texttt{subjects}, e nessun campo li ha tutti e due. Il controllo lo conferma:
riscritta la query con un \texttt{multi\_match} per parola, con gli stessi refusi,
la stessa prima lettera esatta, gli stessi pesi e tutte le parole ancora
obbligatorie, cade solo il vincolo del campo unico, e l'atto giusto torna in
tutte e 300 le query. Koskidex innestato, con le impostazioni di compatibilità,
trova gli stessi insiemi di Elasticsearch su tutte e 300: la parità eredita
anche questo difetto.

\paragraph{Due correzioni, perché i difetti sono due.} Spento in Koskidex
innestato il solo vincolo del campo unico
(\texttt{KOSKIDEX\_ALL\_TERMS\_IN\_ONE\_FIELD=false}), l'atto torna entro i primi
dieci nel 93\% delle query e nessuna resta a vuoto, ma è primo solo nel 38\%.
Nelle 185 query in cui non è primo, in 154 lo scavalca un atto con il numero a un
refuso: per \texttt{1209 Crispiano} vince un atto con \texttt{1109}. È il difetto
della sezione~\ref{sec:numero-atto}, che il vincolo del campo unico nascondeva.
La seconda correzione è un'impostazione di Koskidex con lo stesso significato
di \texttt{disableOnNumbers} di Meilisearch~\cite{meili-typo},
\texttt{disable\_on\_numbers}: un termine della query fatto di sole cifre non
ammette refusi, mentre i codici misti (\texttt{a651}) li tengono. In Documentale
la accende \texttt{KOSKIDEX\_TYPOS\_ON\_NUMBERS=false}.

\begin{table}[ht]
  \centering
  \small
  \begin{tabular}{lrrrr}
    \toprule
    Dall'app & MRR@10 & primo & entro 10 & vuote \\
    \midrule
    Elasticsearch, produzione & 0,017 & 1,3\% & 2,0\% & 70\% \\
    Koskidex, campo unico & 0,018 & 1,7\% & 2,0\% & 70\% \\
    Koskidex, campo unico, n.\ esatti & 0,020 & 2,0\% & 2,0\% & 95\% \\
    Koskidex, parole libere & 0,548 & 38,3\% & 93,3\% & 0 \\
    \textbf{Koskidex, parole libere, n.\ esatti} & \textbf{0,950} & \textbf{91,7\%} & \textbf{99,7\%} & \textbf{0} \\
    \midrule
    Elasticsearch, parole libere (controllo) & 0,534 & 36,3\% & 96,0\% & 0 \\
    \bottomrule
  \end{tabular}
  \caption{Le 300 query \texttt{<numero> <comune>}. \emph{Campo unico}: tutte le
  parole nello stesso campo, come in produzione; \emph{parole libere}: ciascuna in
  un campo qualsiasi; senza \emph{numeri esatti} (n.\ esatti) i numeri ammettono refusi. Le
  prime cinque righe passano dall'app; il controllo interroga Elasticsearch con la
  query riscritta parola per parola.}
  \label{tab:numero-comune}
\end{table}

\textbf{Le due correzioni insieme portano l'atto giusto al primo posto in 275
query su 300}, da 4. Nessuna delle due basta da sola: il vincolo del campo unico
decide se l'atto si trova, i refusi sui numeri se è primo. La riga con il solo
campo unico e i numeri esatti dice anche un'altra cosa: togliere i refusi ai
numeri, lasciando il vincolo, porta le query a vuoto dal 70\% al 95\%. Quel 30\%
che in produzione rispondeva, rispondeva soprattutto con numeri sbagliati.

Nelle 25 query in cui l'atto giusto resta dietro, in 21 il primo classificato ha
lo stesso numero, esatto, in un campo che pesa di più (il numero dell'atto giusto
sta in \texttt{additional\_data}, che pesa 1); in 3 lo ha esatto in un campo che
non pesa di più; in una sola non lo ha esatto, ed è l'unico codice misto fra i
perdenti. Sulle 24 query del confronto il campo libero allarga 14 insiemi su 24,
di 375 atti in tutto, senza toglierne nessuno; i numeri esatti cambiano poi
soltanto le due query con un numero d'atto: \texttt{ordinanza 187} e
\texttt{determina 1223} restituiscono l'atto giusto, primo e unico.

\section{Le date, in tre formati}\label{sec:date}

Nelle schede degli atti la stessa data compare in tre formati, tutti frequenti:
\texttt{14/01/2026} (3.097 volte, in 1.465 atti), \texttt{14 gennaio 2026}
(1.834 volte, in 949) e \texttt{14.01.2026} (2.103 volte, in 982). Gli importi
hanno quasi sempre i punti delle migliaia (312 volte, in 201 atti), e quasi mai
no (2 volte, in 2). Il tokenizer standard di Elasticsearch tiene
\texttt{14.01.2026} come un termine solo e spezza \texttt{14/01/2026} in tre, e
nessun tokenizer sa che \emph{gennaio} è \texttt{01}: chi scrive una data in un
formato non trova gli atti che la scrivono in un altro. Il tokenizer di
Koskidex spezza tutto, e collega barra e punto, ma il giorno di ogni data
diventa un numero che combacia con qualunque ricerca per numero.

\paragraph{Come si misura, e perché diversamente.} Qui il metodo del capitolo
non si può seguire: una query \texttt{<data> <comune>} su Elasticsearch di
produzione andrebbe quasi sempre a vuoto per il difetto della
sezione~\ref{sec:numero-comune}, perché il comune non sta nel campo della data,
e misurerebbe quello. Si misura quindi Koskidex piatto con il tokenizer
standard, che è quello di Elasticsearch, e con il suo; per la parità del
capitolo~\ref{chapter:sistemi} il primo si comporta come Elasticsearch sui
termini. Le query sono sintetiche, come le known-item automatiche: da un atto
si prende una data che in quel comune compare in quell'atto solo, e la si
scrive nei tre formati, cinquanta atti per formato di partenza; allo stesso modo
gli importi, con e senza i punti delle migliaia, su 52 atti: 50 con i punti e i
soli 2 che non li hanno.

\paragraph{Il difetto.} Con il tokenizer standard la data nel formato dell'atto
trova l'atto entro i primi dieci nel 98-100\% delle query; in un altro formato
nello 0-8\%, quasi sempre a vuoto. Con il tokenizer di Koskidex barra e punto si
trovano fra loro (94\%), il mese in lettere no (0-8\%). Gli importi con e senza
punti non si trovano mai.

\paragraph{La correzione.} Due impostazioni di Koskidex, spente per default,
riscrivono il testo degli atti e delle query prima del tokenizer: ogni data
valida, in uno dei quattro formati e con i mesi in italiano, diventa un numero
solo, \texttt{aaaammgg}, e agli importi si tolgono i punti delle migliaia. Con
tutti e due i tokenizer ogni coppia di formati trova l'atto entro i primi dieci
nel 98-100\% dei casi, gli importi nel 100\%, e le 300 ricerche
\texttt{<numero> <comune>} non peggiorano (MRR@10 da 0,975 a 0,977). Due
previsioni su otto sono cadute. I numeri piccoli, da 1 a 31, perdono meno
candidati del previsto, perché il comune nella query li restringe già. E
l'indicizzazione costa il 29-34\% in più: avevo previsto al più il 20\%, e
saltare i testi che non possono contenere un formato l'ha solo ridotta dal
41-46\%, perché quasi ogni scheda contiene davvero cifre, barre e punti.

Le query sono fabbricate: dicono che il motore trova una data in qualunque
formato, non quanto spesso una persona cerca per data, e nelle known-item umane
raccolte finora non succede mai: nessuna delle 76 contiene una data. E la correzione ha un prezzo che nessuna
previsione misurava: la data diventa un termine solo, quindi il suo anno non
combacia più con una ricerca che contiene l'anno da solo, come \emph{bilancio
2026}, a meno che l'anno compaia anche fuori da una data.

\paragraph{Dall'app.} In Documentale le due impostazioni sono due variabili,
spente e fuori dal profilo della sezione seguente, misurate dall'app con il
profilo acceso (\texttt{2026-09-28\_date-app}). Le date si comportano come in
Koskidex piatto: da 0-8\% a 98-100\% di atti entro i primi dieci in un altro
formato, MRR@10 da 0,315 a 0,979, con il 6\% in più di indicizzazione; le
known-item automatiche migliorano appena, le 76 umane raccolte finora e le 21
query del confronto senza cifre non cambiano di un posto. Gli importi no, e
per una ragione che Koskidex piatto non poteva mostrare, perché
\texttt{scripts/evaluate} cerca senza refusi e l'app con la tolleranza
automatica: nell'app gli importi si trovano fra formati anche spenti,
perché \texttt{5056,03} dista un carattere da \texttt{5.056,03}, e un termine
con la virgola non è fatto di sole cifre, quindi i numeri esatti del
profilo non lo riguardano. Lo stesso refuso fa combaciare importi diversi, e
togliere i punti peggiora le cose: gli importi diventano più corti e più
vicini, e l'MRR@10 degli importi scende da 0,917 a 0,891. Il difetto vero
degli importi è che ammettono refusi.

\paragraph{Gli importi senza refusi.} È il difetto della
sezione~\ref{sec:numero-atto} su un'altra forma di numero, e la correzione è la
stessa: un'impostazione di Koskidex, \texttt{disable\_on\_amounts}, spenta per
default, nega i refusi ai termini fatti di cifre, punti e virgole, e in
Documentale è una variabile fuori dal profilo
(\texttt{2026-09-28\_importi-esatti}). Da sola toglie il ponte che il refuso
dava per caso, e gli importi fra formati non si trovano più (0 su 52); insieme
alla normalizzazione degli importi si trovano tutti (52 su 52), l'MRR@10 degli
importi sale da 0,917 a 0,990, e i risultati sulle 104 query scendono da 175 a
106: quelli tolti sono atti con un importo diverso. Sette previsioni su otto;
quella caduta è caduta in meglio. Anche \texttt{14.01.2026} è un termine con i
punti, che prima ammetteva refusi e trovava \texttt{14.01.2025}: con
l'impostazione 14 ricerche di date col punto migliorano, nessuna peggiora, e i
risultati sulle query delle date scendono da 594 a 241. Le known-item
automatiche, le umane raccolte finora e le 21 query del confronto senza cifre
non cambiano.

Nell'app quindi le date vanno normalizzate, e gli importi normalizzati ed
esatti. Del prezzo dell'anno da solo le query raccolte finora non dicono niente:
nessuna delle umane contiene un anno, e l'unica del confronto che lo contiene
non trova niente in nessun caso. Per questo le tre impostazioni restano fuori
dal profilo finché non girano le known-item umane complete.

\section{Cosa cambia per chi cerca}\label{sec:cosa-cambia}

Le tre correzioni misurate dall'app sono misurate una alla volta, e una configurazione fatta di
pezzi misurati separatamente potrebbe non fare la somma dei pezzi. In Documentale
esiste per questo un profilo, \texttt{KOSKIDEX\_PROFILO=consigliata}, che fa
partire insieme le tre variabili dal valore corretto, senza cambiare nessun
default: senza profilo l'app si comporta come in produzione, e ogni variabile
impostata da sola vince sul profilo. Il profilo è misurato come un tutto, da un
indice vuoto, con le stesse query:
\begin{itemize}
  \item sulle 300 query known-item fa \textbf{esattamente} quello che facevano
        parole libere e numeri esatti, fino all'ordine dei primi dieci risultati di
        ogni query: l'elisione non sposta niente;
  \item sulle 24 query del confronto, rispetto a parole libere e numeri esatti,
        nessuna query perde un documento e 7 ne guadagnano, quelle toccate
        dalle elisioni (l'ottava della sezione~\ref{sec:elisioni},
        \texttt{ordinanza 187}, con i numeri esatti ha già il solo atto giusto);
  \item \texttt{ordinanza 187} e \texttt{determina 1223} danno il solo atto giusto.
\end{itemize}

\begin{table}[ht]
  \centering
  \begin{tabular}{lrr}
    \toprule
    & produzione & consigliata \\
    \midrule
    \texttt{ordinanza 187} & 12° su 85 & 1° su 1 \\
    \texttt{determina 1223} & 5° su 8 & 1° su 1 \\
    \texttt{<numero> <comune>}, atto primo & 1,3\% & 91,7\% \\
    \texttt{<numero> <comune>}, a vuoto & 70\% & 0 \\
    coppie (atto, parola) perse per l'elisione & 5.775 & 36 \\
    atti con una parola introvabile & 40,3\% & 0,4\% \\
    \bottomrule
  \end{tabular}
  \caption{La ricerca di Documentale in produzione e con il profilo consigliato di
  Koskidex innestato. Le righe delle elisioni vengono dalla misura delle 507
  parole della sezione~\ref{sec:elisioni}.}
  \label{tab:cosa-cambia}
\end{table}

La tabella~\ref{tab:cosa-cambia} riassume cosa cambia. Il numero degli atti con una
parola introvabile, con il profilo, è quello delle 36 coppie residue, che stanno in
36 atti diversi.

\paragraph{Cosa non dice.} Tutte queste misure riguardano \emph{se} l'atto si
trova e, per le query identificative, \emph{dove}. Non dicono niente dell'ordine
dei risultati sulle ricerche in lingua naturale, come \emph{manutenzione strade}:
lì non c'è un atto giusto solo, e per misurare servono giudizi di pertinenza, che
sono l'oggetto del capitolo~\ref{chapter:confronto}. Non dicono neanche che
Elasticsearch non si possa correggere: il filtro delle elisioni e la query
riscritta parola per parola sono misurati anche lì, con lo stesso effetto. Per i
numeri servirebbe separare nella query i termini fatti di sole cifre e cercarli
senza refusi, e questo in Elasticsearch non è stato misurato.

\paragraph{Cosa vuol dire.} I quattro difetti hanno una cosa in comune: nessuno è
un errore di programmazione. Il recupero congiuntivo per campo, la tolleranza ai
refusi uguale per parole e numeri, il tokenizer che rispetta l'apostrofo e tiene
insieme le date con il punto sono tutte impostazioni ragionevoli di un motore
generico, e ciascuna si rompe su una proprietà precisa dei documenti
amministrativi italiani: i metadati sparsi in campi diversi, i numeri di
protocollo, gli articoli elisi, le date scritte in tre modi. Per trovarle è bastato
fare le ricerche che fa chi usa un documentale e contare, con il codice dell'app,
quante andavano a vuoto.
